离线指标好、上线之后一路下滑，而且查不出任何自然原因——这种情形通常意味着对面有人在读你的输出并据此调整。这一刻，统计学习的整套语言会同时失效：没有固定的分布，就没有真风险；没有真风险，泛化界断言的那个量根本不存在。问题不在于界不够紧，在于它谈论的对象消失了。

本单元换一个能在这种环境里兑现的承诺。它不问``我们离真值多远''，只问``事后回头看，早知道该一直用哪个决策，我们比它多亏了多少''。这个量叫遗憾。它对世界不作任何假设，所以即使对手完全恶意也算数；相应地它也弱得多，因为参照系被换成了一个我们自己选定的基准。整个单元就在这条交换线上展开：承诺换了参照系，能保证的东西随之改变，而边界必须逐条说清。

\hyperref[sec:17-Constrained-Guarantees/04-Online-Learning-and-Regret/01]{``对手给数据时该比什么''}说明 i.\,i.\,d.\ 假设为什么是第一个塌掉的东西，并把遗憾这个替代承诺的四条弱点摆开。\hyperref[sec:17-Constrained-Guarantees/04-Online-Learning-and-Regret/02]{``无遗憾是可以做到的''}给出两个具体算法——指数权重与在线梯度下降——各自的证明骨架与成立条件，并证明 \(\sqrt T\) 是任何算法都绕不过去的下界。\hyperref[sec:17-Constrained-Guarantees/04-Online-Learning-and-Regret/03]{``探索与利用是同一笔预算''}掐掉一条信息通道：只看得见自己选中那个动作的损失时，最优遗憾从 \(\sqrt{T\ln N}\) 变成 \(\sqrt{TN}\)，信息的缺失被换算成了一个可以写出来的数。

最后一节结束在一个新问题上：如果对手不是恶意的，而是另一个同样在学习的人，两套无遗憾算法互相对打会走向哪里。下一个单元接手这件事。
